\chapter{Enhanced Methodology and System Architecture: AgentClinic v2}
\label{Chapter5}
\lhead{Chapter 5. \emph{Enhanced Methodology \& System Architecture (v2)}}

To address the architectural flaws identified in the initial framework, we developed \textbf{AgentClinic v2} as a complete redesign of the system.

\section{LangGraph Directed Cyclic Graph (DCG) Orchestration}

The core limitation of the prior ad-hoc loop was its inability to enforce strict, deterministic state transitions. In v2, the orchestrator was rebuilt as a \textbf{LangGraph Directed Cyclic Graph (DCG)}, where the consultation protocol is enforced as a finite-state controller. Nodes represent agents, and edges encode deterministic routing logic conditioned on the current shared state.

\subsection{The Latent Metacognitive Reflection Node}

The most important architectural change is the introduction of the \texttt{doctor\_reflection} node. Human clinicians maintain a mental differential diagnosis that evolves silently as they interview a patient. Standard LLMs have no equivalent scratchpad; when forced to generate bedside dialogue and medical reasoning simultaneously, they tend to hallucinate or leak internal reasoning into spoken dialogue.

In v2, the DCG routes execution to the \texttt{doctor\_reflection} node before any spoken utterance. Here, the LLM outputs a structured JSON array of its current top-3 differential diagnoses (e.g., \texttt{["Asthma", "COPD", "Pneumonia"]}), operating at a low sampling temperature ($T=0.1$) to ensure deterministic extraction. These hypotheses are stored in the \texttt{ClinicalState} trajectory log but are \textit{not} injected into the patient-facing dialogue history. This makes the LLM's latent reasoning observable and quantifiable without contaminating the interaction.

\begin{figure}[H]
\centering
\begin{tikzpicture}[
    node distance=2.2cm and 3.2cm,
    agent/.style={rectangle, rounded corners, draw=black, thick,
                  minimum width=3.2cm, minimum height=1cm, align=center, font=\small},
    hidden/.style={rectangle, rounded corners, draw=black, fill=gray!20, thick,
                   minimum width=3.2cm, minimum height=1cm, align=center, font=\small},
    terminal/.style={rectangle, rounded corners, draw=black, fill=gray!40, thick,
                     minimum width=2.2cm, minimum height=1cm, align=center, font=\small},
    arrow/.style={thick, ->, >=stealth}
]
\node[hidden] (reflection) {\textbf{doctor\_reflection} \\ (Latent DDx)};
\node[agent, right=of reflection] (speaker) {\textbf{doctor\_speaker} \\ (Utterance)};
\node[agent, below left=of speaker] (patient) {\textbf{patient\_node}};
\node[agent, below right=of reflection] (measurement) {\textbf{measurement\_node}};
\node[terminal, right=of speaker] (end) {END};

\draw[arrow] (reflection) -- node[above, font=\small]{Updates State} (speaker);
\draw[arrow] (speaker) -- node[above right, sloped, font=\small]{REQUEST TEST} (measurement);
\draw[arrow] (speaker) -- node[above left, sloped, font=\small]{Ask Question} (patient);
\draw[arrow] (speaker) -- node[above, font=\small]{DIAGNOSIS READY} (end);
\draw[arrow] (speaker.north) .. controls +(up:1cm) and +(up:1cm) ..
    node[above, font=\small]{turn $\geq$ max} (end.north);
\draw[arrow] (patient) -- node[left, font=\small]{Response} (reflection);
\draw[arrow] (measurement) -- node[left, font=\small]{Results} (reflection);
\end{tikzpicture}
\caption{LangGraph DCG topology of AgentClinic v2. The \texttt{doctor\_reflection} node (shaded) ensures reasoning trajectories are recorded systematically before each spoken turn, providing a structured metacognitive scaffold without contaminating patient-facing dialogue.}
\label{fig:langgraph}
\end{figure}

\section{Heterogeneous Multi-Engine Architecture and Quantisation}

To eliminate \textit{Lexical Leakage}, v2 follows a \textit{Cognitive Asymmetry} design principle: the Doctor operates as a domain-expert engine, while the Patient, Measurement, and Moderator agents run on separate, architecturally distinct model instances.

Running three or more distinct 7B-parameter models simultaneously on a single NVIDIA H100 NVL (95.8\,GB VRAM) is impractical in standard FP16, since each model requires roughly 14\,GB excluding activations. We therefore adopted \textbf{NormalFloat4 (NF4) quantisation} via the BitsAndBytes library. Since neural network weights approximate a zero-mean normal distribution, NF4 maps them into an information-theoretically optimal 4-bit representation. On top of this, \textit{double quantisation} applies a second round of quantisation to the quantisation constants themselves, recovering an additional $\approx$0.37 bits per parameter.

\begin{equation}
\text{VRAM}_{\text{model}} \approx \text{Params} \times \left( \frac{4\ \text{bits}}{8\ \text{bits/byte}} \right) + \text{KV-Cache Overhead}
\end{equation}

This reduces the memory footprint per 7B model by approximately 86.4\% (from $\approx$28\,GB in FP32 to $\approx$3.8\,GB), making heterogeneous multi-engine isolation feasible and leaving room for Flash Attention 2 to handle long-context processing.

\section{Formulation of Process-Oriented Clinical Metrics}

Binary diagnostic accuracy alone can mask serious reasoning flaws. The \texttt{ClinicalTrajectoryEvaluator} computes three process-oriented metrics inline, using data accumulated in the \texttt{ClinicalState} trajectory.

\paragraph{1. Diagnostic Stability (DS):} Measures the average Jaccard similarity between consecutive differential diagnoses over $T$ turns. High stability suggests logical hypothesis refinement; low stability signals erratic switching.
\begin{equation}
DS = 1 - \frac{1}{T - 1} \sum_{t=1}^{T-1} \left( 1 - \frac{|DDx_t \cap DDx_{t+1}|}{|DDx_t \cup DDx_{t+1}|} \right)
\end{equation}

\paragraph{2. Test Rationality (TR):} A binary compliance indicator that checks whether the model ordered at least one basic investigation (e.g., CBC, vitals) before requesting expensive or invasive procedures (e.g., MRI, biopsy). Let $t_{\text{basic}}$ and $t_{\text{adv}}$ represent basic and advanced tests, respectively:
\begin{equation}
TR = \begin{cases}
0 & \text{if } \exists\, t_{\text{adv}} \text{ ordered before any } t_{\text{basic}} \\
1 & \text{otherwise}
\end{cases}
\end{equation}

\paragraph{3. Information Efficiency (IE):} The ratio of unique differential hypotheses generated ($N_{\text{unique}}$) to the total dialogue turns utilised ($T$):
\begin{equation}
IE = \frac{N_{\text{unique}}}{T}
\end{equation}

\section{Structured Cognitive Bias Injection Manifold}

The v2 framework carries forward and enhances the bias injections from the initial framework. Instead of generic sentence-level prompts, we inject ecologically valid distortions:

\begin{itemize}
    \item \textbf{Recency Bias:} The system maintains a running queue of the $k=5$ most recently processed case summaries. These are injected into the Doctor's system prompt, mimicking a physician overly influenced by recent clinical encounters. This exploits the transformer's natural tendency to assign elevated attention weights to tokens near the current generation boundary.
    \item \textbf{Confirmation Bias:} The Doctor is primed with a strong, incorrect initial hypothesis (e.g., \textit{``You are initially confident the patient has malignancy.''}). This warps the model's forward probability distribution toward confirmatory evidence acquisition, operationalising the premature diagnostic closure seen in anchoring studies.
\end{itemize}

\section{Task-Specific Dynamic LoRA Adapter Routing}

To address the catastrophic forgetting observed in narrowly fine-tuned models, we implemented a parameter-efficient inference-time solution: \textbf{dynamic routing of task-specific Low-Rank Adaptation (LoRA) adapters}, conditioned on the active LangGraph node.

LoRA~\cite{hu2022lora} decomposes a weight update $\Delta W \in \mathbb{R}^{d \times k}$ into a low-rank product $\Delta W = BA$, where $B \in \mathbb{R}^{d \times r}$ and $A \in \mathbb{R}^{r \times k}$ with rank $r \ll \min(d, k)$. The effective update applied to the frozen pre-trained matrix $W_0$ is scaled by $\frac{\alpha}{r}$:
\begin{equation}
W = W_0 + \frac{\alpha}{r} \cdot BA
\end{equation}
where $\alpha = 32$ and $r = 16$, yielding a scaling factor of $2\times$. This reduces the number of trainable parameters by orders of magnitude while preserving representational capacity in the task-relevant subspace. LoRA adapters were applied to all four attention projection matrices (\texttt{q\_proj}, \texttt{k\_proj}, \texttt{v\_proj}, \texttt{o\_proj}) across all 32 transformer layers, totalling approximately 16.8M trainable parameters ($\sim$25\,MB) per adapter.

Two lightweight adapters were trained independently using QLoRA (4-bit NF4 quantised base with BF16 adapter training) via the HuggingFace PEFT library:
\begin{itemize}
    \item A \textit{reasoning adapter} ($r=16$, $\alpha=32$), fine-tuned on medical differential diagnosis corpora and activated exclusively during the \texttt{doctor\_reflection} node via \texttt{model.set\_adapter("reasoning")}. This adapter steers the model's latent probability distribution toward structured diagnostic enumeration at a low temperature ($T = 0.1$).
    \item A \textit{dialogue adapter} ($r=16$, $\alpha=32$), fine-tuned on empathetic clinical communication transcripts and activated during the \texttt{doctor\_speaker} node via \texttt{model.set\_adapter("dialogue")}. This adapter preserves the conversational fluency and empathetic tone needed for realistic patient interaction at standard temperature ($T = 0.7$).
\end{itemize}

\textbf{Data separation for experimental validity:} To prevent data leakage (train-test contamination), we trained the adapters on clinical cases from \texttt{agentclinic\_nejm.jsonl} (New England Journal of Medicine scenarios), while all evaluation experiments (E1--E6) used \texttt{agentclinic\_medqa.jsonl} (107 MedQA scenarios) exclusively. This strict separation ensures that the adapter's contribution to E6 accuracy (54.7\%) reflects genuinely transferable clinical competence rather than memorisation of evaluation answers.

At inference time, the DCG invokes \texttt{model.set\_adapter("reasoning")} before reflection and \texttt{model.set\_adapter("dialogue")} before speech generation. The adapter swap is a pointer operation ($<$1\,ms latency) that redirects which $BA$ matrices are added to the frozen $W_0$. The combined VRAM overhead of both adapters is approximately 50\,MB, which is negligible relative to the base model footprint. This allows the dual-adapter paradigm to operate within the same memory envelope as a single unmodified model.

\section{Theoretical Underpinnings of the Architecture}

\subsection{Decoupling Reasoning from Generation (Latent Chain-of-Thought)}

Standard auto-regressive LLMs compute the probability of each token conditioned on all preceding tokens: $P(w_t \mid w_1, \ldots, w_{t-1})$. When a model is asked to simultaneously generate bedside dialogue and perform complex multi-conditional reasoning, its self-attention heads must split capacity between syntactic fluency and semantic medical deduction. The result is often reasoning-contaminated utterances or outright hallucinations.

The \texttt{doctor\_reflection} node acts as a \textit{Latent Chain-of-Thought (CoT) scaffold}~\cite{wang2024survey}. By first forcing the LLM to output a JSON differential, high-information medical tokens are injected directly into the context window. When the \texttt{doctor\_speaker} node subsequently executes, its attention heads attend heavily to this structured JSON, constraining the dialogue distribution to align with the prior diagnostic reasoning.

\subsection{Embedding Manifold Overlap and Lexical Leakage}

LLMs project text into high-dimensional latent vector spaces. In a homogeneous multi-agent setup, the Doctor and Patient share the exact same model weights ($\theta_{\text{shared}}$), so their latent representations of symptoms, anatomical structures, and diseases occupy identical embedding manifolds.

\textit{Lexical Leakage} arises because the Doctor agent can perform zero-shot knowledge transfer by recognising the mathematical proximity of the Patient's token embeddings within this shared manifold. The heterogeneous multi-engine architecture severs this overlap. By forcing the Doctor to interact with a Patient agent built on a completely different neural architecture, we ensure that the Doctor must perform genuine semantic deduction rather than exploiting implicit geometric similarity.

\subsection{Information-Theoretic Basis of NF4 Quantisation}

Pre-trained neural network weights typically follow a zero-mean normal (Gaussian) distribution. Standard uniform quantisation (such as INT8) wastes bit capacity on sparse outlier regions of this distribution. NF4 creates non-linearly spaced quantisation bins that are information-theoretically optimal for normally distributed data, placing more bins in the high-density central region where the majority of weights reside. This preserves the statistical integrity of the LLM's clinical reasoning while dramatically reducing VRAM footprint.
